\documentclass[10pt]{article}
\usepackage{tmlr}
\usepackage{amsmath}
\usepackage{booktabs}
\usepackage{hyperref}
\usepackage{url}
\hypersetup{hidelinks}

\title{Sentence Boundaries Are Part of the Watermark Interface:
A SWAN Audit and an Exact Resegmentation Framework}

\author{\name Anonymous Authors}

\def\month{10}
\def\year{2026}
\def\openreview{\url{https://openreview.net/forum?id=XXXX}}

\begin{document}

\maketitle

\begin{abstract}
Sentence-level semantic watermarks treat a segmentation of text as detector
input, yet robustness is usually reported for a single segmentation. We study
this interface at behavioral, algorithmic, and systems levels. First, we audit
SWAN. On 30 saved five-unit texts, replacing four sentence-final punctuation
marks with semicolons while preserving lexical tokens and order changed the
number of below-threshold decisions from 1/30 to 30/30; restoring the saved
five spans returned it to 1/30. This establishes a reversible short-text
interface dependence, but SWAN's one-unit score ceiling already explains the
fixed-threshold outcome, and two normal-length count-neutral routes stopped
before detection. Second, we formalize fixed-token, fixed-unit-count boundary
relocation and give exact dynamic programs for independent-span and
first-order adjacent-span detectors. The pairwise program agrees with brute
force on 7,669 systematic and random instances, agrees with the additive
program on 1,000 reductions, and solves a 300-token case containing
14,606,509 legal segmentations in 0.8054 seconds. A compiled factorized
implementation exactly processes a synthetic real-data shape representing
$3.94\times10^{19}$ segmentations. Third, we stress-test the full path on
SemStamp and SimMark. A final frozen SemStamp detector-health run obtained
4.7\% test FPR and 95.4\% positive TPR, but exact-certificate readiness
stopped before scoring because one document required 78.6 million pruned
span pairs and the two-replay budget was 86.65 GPU-hours. We therefore report
no real certified rate. Together, the results show why segmentation must be
specified as part of the detector contract, what exact certification can
guarantee, and where an audited implementation still becomes impractical.
\end{abstract}

\section{Introduction}

Semantic text watermarks aim to preserve detectability under lexical or
syntactic variation by placing signals in sentence-level representations.
This design introduces an interface that token-level evaluations can obscure:
before a sentence-level detector can score a document, some procedure must
decide what its units are. The same word sequence may consequently induce
different statistical trials, different encoder or parser calls, and
different pairwise relations.

We investigate this problem through one research program with two evidence
streams. The first is an empirical audit of SWAN, an AMR-based semantic
watermark \citep{ye2026swan}. The second asks whether robustness to a declared
set of word-preserving resegmentations can be certified exactly for
sentence-level detectors such as SemStamp \citep{hou2024semstamp},
k-SemStamp \citep{hou2024ksemstamp}, and SimMark
\citep{dabiriaghdam2025simmark}. These streams address the same interface at
different levels, but their evidence is not interchangeable: a SemStamp
certificate does not establish a SWAN mechanism, and a SWAN counterexample
does not establish that a certificate is non-vacuous.

The SWAN audit begins with a reversible behavior observation. Thirty saved
texts contain five detection units each. Replacing four sentence-final marks
with semicolons preserves lexical tokens and their order while making the
default detector see one unit. Below-threshold decisions change from 1/30 to
30/30. Supplying the detector with the five saved spans restores the count to
1/30. However, with SWAN's default null rate, a one-unit document cannot
cross the strict public threshold even if its unit is marked. This is a real
interface failure, not evidence that AMR semantics are uniquely fragile.

We attempted to isolate effects beyond unit-count reduction through
normal-length, count-neutral boundary exchange. Two preregistered generation
routes stopped before transformation and detection. A later targeted
construction produced one detectable 126-word document, but not the
source-anchored, independently reviewed cohort required for confirmation. The
normal-length causal endpoint therefore remains unestimated.

The certification stream turns the same concern into a finite optimization
problem. Given a frozen token tape, original boundaries, admissible spans, a
fixed number of units, and a relocation radius $k$, we minimize the detector
score over every legal segmentation in that set. Independent-span scores
admit a standard segmentation dynamic program. SemStamp, k-SemStamp, and
SimMark instead use first-order relations between adjacent units, requiring
an edge-state dynamic program that retains the preceding span.

We validate the algorithms against independent exhaustive oracles, implement
split-conformal calibration with strict ties, reproduce public detector
interfaces on frozen assets, and carry the pipeline to its final resource
gate. The terminal SemStamp run is informative precisely because it does not
produce a certificate: detector health passes, while exact workload capacity
and budget do not. This separates ``the detector is unhealthy,'' ``the exact
algorithm is wrong,'' and ``the validated implementation is too expensive''
as distinct failure modes.

Our contributions are:

\begin{enumerate}
  \item A scope-limited SWAN audit showing a reversible sentence-boundary
  dependence on 30 short texts, together with the exact one-unit statistical
  ceiling and complete normal-length stopping provenance.
  \item A fixed-token, fixed-unit-count boundary-relocation threat model and
  exact dynamic programs for independent-span and first-order adjacent-span
  detectors, including a direct soundness argument.
  \item Oracle, reduction, calibration, adapter, and compiled-core validation
  that distinguishes mathematical correctness from real-model cost.
  \item An end-to-end reproduction record for SemStamp and SimMark, including
  detector-health results and a score-blind SemStamp readiness Stop. We do
  not claim a real certified rate.
\end{enumerate}

\section{Related Work}

\paragraph{Sentence-level semantic watermarks.}
SemStamp partitions a sentence-embedding space with locality-sensitive
hashing and conditions each accepted region on the preceding sentence
\citep{hou2024semstamp}. k-SemStamp replaces that partition with clustering
\citep{hou2024ksemstamp}. SimMark embeds a signal through adjacent-sentence
similarities and uses a soft detector \citep{dabiriaghdam2025simmark}. SWAN
injects AMR templates and detects them through sentence segmentation, neural
AMR parsing, graph normalization, S2Match, and document aggregation
\citep{ye2026swan}.

\paragraph{Resegmentation and rewriting.}
AliMark explicitly addresses sentence splitting and merging through
restructuring and bit-sequence alignment \citep{li2026alimark}. SwordStamp is
the closest structural neighbor: it evaluates rewording, reordering, and
resegmentation, studies unit-count effects, and proposes order-robust
sub-sentence units \citep{diaa2026swordstamp}. These works preclude claims
that we introduce merge/split attacks, boundary exchange, or
segmentation-aware detection. Our distinction is exhaustive minimization over
a declared finite resegmentation set rather than a sampled or heuristic
attack set.

DIPPER provides controllable paragraph paraphrasing
\citep{krishna2023dipper}; SIRA targets high-self-information tokens
\citep{cheng2025sira}; RLCracker studies adaptive worst-case watermark
vulnerability \citep{huang2026rlcracker}; and Chainwash studies repeated
rewriting for diffusion-language-model watermarks
\citep{ameen2026chainwash}. WaterPark emphasizes joint reporting of
detection, quality, and realized false-positive rates
\citep{liang2025waterpark}. FreqMark and CoheMark provide further
sentence-level context \citep{xu2024freqmark,zhang2025cohemark}.

\paragraph{Calibration.}
We use split-conformal order statistics to freeze thresholds before
independent testing \citep{angelopoulos2023gentle}. Binomial uncertainty is
reported with one-sided Clopper--Pearson bounds
\citep{clopper1934confidence}. Calibration is part of the certificate
procedure: the worst-case score distribution is not silently assigned the
raw detector threshold when comparing power at a common target FPR.

\section{Study Design and Evidence Separation}

\subsection{One question, two evidence streams}

The unifying question is whether sentence boundaries are an implicit part of
a sentence-level watermark's security contract. The empirical stream asks
whether changing that interface changes a deployed decision. The
certification stream asks whether a detector can guarantee its decision for
every segmentation in a declared neighborhood.

\begin{table}[t]
\caption{Evidence roles in the unified study.}
\label{tab:evidence}
\centering
\small
\begin{tabular}{p{0.20\linewidth}p{0.34\linewidth}p{0.37\linewidth}}
\toprule
Tier & Evidence & Permitted use \\
\midrule
SWAN behavior & 30 short-text punctuation A--B--A pairs &
System-specific interface result at a historical threshold \\
SWAN stopping & Two eight-source normal-length pilots &
Feasibility and stopping provenance, not detector evidence \\
Exact algorithms & Oracle and reduction tests on frozen lattices &
Correctness for the stated mathematical detector classes \\
Synthetic statistics & 99,980 synthetic null documents &
Calibration implementation, not real-system FPR \\
System reproduction & Frozen public SemStamp/SimMark assets &
Implementation equivalence and detector-health diagnosis \\
Certificate readiness & 3,000-row SemStamp workload audit &
Capacity and cost conclusion, not a certified robustness rate \\
\bottomrule
\end{tabular}
\end{table}

The streams appear in one submission because the audit motivates the formal
problem and the formalization clarifies what the audit did not establish.
They remain separate at the claim level. In particular, no SemStamp or
SimMark result upgrades the missing SWAN count-neutral endpoint.

\subsection{Fail-closed protocol}

Each stage freezes inputs, code identity, thresholds, continuation gates, and
resource limits before reading the next class of result. A failed gate
terminates that version. We do not delete hard samples, relax a radius,
increase a budget, or substitute a cheaper bound while preserving the old
version label. This produces many negative engineering records, but prevents
partial runs from becoming post-hoc positive claims.

\section{Part I: Auditing the SWAN Segmentation Interface}

\subsection{Detection chain and statistical boundary}

Let a segmenter map a document to units $s_1,\ldots,s_N$. We distinguish:

\begin{itemize}
  \item $B$, SWAN's fixed secret bank of AMR templates;
  \item $R(s)$, the frozen parser's raw AMR output for unit $s$;
  \item $T(R(s))$, the normalized candidate graph supplied to S2Match;
  \item $D(g_1,g_2)$, an artificial composition of saved graphs used only for
  scorer diagnostics.
\end{itemize}

A unit is \emph{green} when at least one fixed-bank match crosses the unit
threshold. If $G$ of $N$ units are green, the public statistic is

\begin{equation}
z = \frac{G-\lambda N}{\sqrt{N\lambda(1-\lambda)}}.
\label{eq:swan}
\end{equation}

For fixed $N$,

\begin{equation}
z_{\max} = \sqrt{\frac{N(1-\lambda)}{\lambda}}.
\end{equation}

At $\lambda=0.25$, $z_{\max}=\sqrt{3}=1.7321$ for $N=1$, below the
strict threshold $z>2.33$. SWAN also has two entry-point semantics: the
direct-AMR path returns \texttt{NaN} for $N<2$, while the pre-parsed path
evaluates Equation~\ref{eq:swan}. We preserve that distinction and never
count \texttt{NaN} as an attack success.

\subsection{Short-text A--B--A intervention}

The saved cohort contains 30 independently generated SWAN baselines of
17--46 words and five detection units. The intervention replaces four
sentence-final punctuation marks with semicolons without changing lexical
tokens or their order. We compare the original text, the punctuation-merged
text under default segmentation, and the same merged text supplied with the
five saved original spans.

\begin{table}[t]
\caption{Saved SWAN A--B--A behavior at the strict public threshold.}
\label{tab:swan-aba}
\centering
\begin{tabular}{lrr}
\toprule
Detector input & Units & $z\leq2.33$ \\
\midrule
Original text & 5 & 1/30 \\
Terminal marks replaced by semicolons & 1 & 30/30 \\
Merged text with saved spans restored & 5 & 1/30 \\
\bottomrule
\end{tabular}
\end{table}

The paired changes in both directions have exact $p=3.73\times10^{-9}$.
Table~\ref{tab:swan-aba} establishes that external segmentation can
reversibly determine the final decision for these texts and this entry point.
It does not establish an AMR-specific mechanism: the middle condition changes
$N$ from five to one, and the one-unit bound is already sufficient.

A separate no-feedback three-arm study recorded 29/30 valid attacks in its
boundary-reorganization arm, but those generated candidates could contain
lexical changes. Historical fixed-budget and feedback-assisted studies
recorded 191/250 and 218/250 valid attacks, respectively. They motivate the
audit but do not replace the controlled A--B--A result.

\subsection{Normal-length count-neutral design and stops}

The confirmatory intervention would merge one adjacent sentence pair and
split another sentence at a detector-blind natural boundary. It preserves
tokens, order, punctuation multiset, and unit count:
$N(C_2)=N(C_0)$. Quality must be frozen by two independent human reviewers
before detection, and each changed detection procedure requires independent
null calibration.

The \emph{final1} pilot assigned eight fresh sources. All produced
100--220-word terminal documents, six passed structural checks, and none of
those six passed the frozen model-assisted quality labels. It stopped before
transformation, parsing, matching, or detection. Those labels are not
reported as human quality estimates.

An independent source-anchored route bound ten generation slots to ten source
anchors per document. Five of eight documents passed structure, below the
six-of-eight continuation gate. All 60 completed slots used fallback, and
three additional slots exhausted 30 candidates. It stopped before quality
review, transformation, and detection.

\begin{table}[t]
\caption{Normal-length SWAN routes and their terminal evidence.}
\label{tab:swan-stops}
\centering
\small
\begin{tabular}{p{0.18\linewidth}rrp{0.43\linewidth}}
\toprule
Route & Structural pass & Detection runs & Terminal interpretation \\
\midrule
final1 & 6/8 & 0 & Auxiliary quality gate failed; no paired cohort \\
Source-anchored & 5/8 & 0 & Structural gate failed before quality review \\
Target-47, 126 words & 1 document & 1 baseline &
20/20 green, but targeted and not a confirmatory cohort \\
\bottomrule
\end{tabular}
\end{table}

A later engineering study showed that one targeted, single-template
construction can be both normal length and detectable: a 126-word,
20-unit text had 20/20 strict injections and 20/20 green units
($z=7.7460$). It lacked source anchoring, independent human review, an
independent cohort, and empirical FPR. Its standard count-neutral
transformation failed because no admissible natural split was available.
An auxiliary equal-$N$ boundary relocation changed 19/19 to 18/19 green,
leaving both versions far above threshold. This closes a narrow engineering
question, not the confirmatory endpoint.

\subsection{SWAN evidence boundary}

The normal-length effect beyond unit-count reduction, the contribution of
real $R(s)$ and $T(R(s))$, and condition-specific FPR remain unestimated.
Artificial graph composition is therefore confined to Appendix~A. The
finding is a system-interface counterexample, not a general watermark-removal
algorithm and not evidence that sentence meaning changed.

\section{Part II: Exact Resegmentation Certificates}

\subsection{Threat set}

Let $x=(x_1,\ldots,x_L)$ be a frozen token tape. An original segmentation
$P_0$ has internal boundary set $B_0$ and fixed unit count
$m=|B_0|+1$. Let $\mathcal G$ be the admissible internal gaps, frozen before
scores are observed. Each candidate span must satisfy a declared legality
predicate, such as a native tokenizer context bound.

For any segmentation $P$ with $|B(P)|=|B_0|$, define the relocation radius

\begin{equation}
\rho(P,P_0)=|B(P)\setminus B_0|
           =\frac{1}{2}|B(P)\triangle B_0|.
\end{equation}

The complete radius-$k$ threat set is

\begin{equation}
\mathcal A_k(x,P_0)=
\left\{P:
\begin{array}{l}
B(P)\subseteq\mathcal G,\quad |B(P)|=|B_0|,\\
\rho(P,P_0)\leq k,\quad
\text{every span in }P\text{ is legal}
\end{array}
\right\}.
\label{eq:threat}
\end{equation}

The threat model preserves token identity, token order, and unit count. A
relocation can move a boundary any distance; it is not a local character-edit
metric. The certificate does not cover paraphrasing, insertion, deletion,
variable $N$, arbitrary punctuation rendering, or a natural-language
segmenter's reachability.

\subsection{Detector classes}

For independent-span scores $s(i,j)$,

\begin{equation}
D(x,P)=b+\frac{1}{m}\sum_{[i,j)\in P}s(i,j).
\end{equation}

For first-order adjacent-span detectors, let
$0=b_0<b_1<\cdots<b_m=L$ be the segmentation endpoints and let
$q(h,i,j)$ score adjacent spans $[h,i)$ and $[i,j)$. Then

\begin{equation}
D(x,P)=b+\frac{1}{c}
\sum_{t=2}^{m}q(b_{t-2},b_{t-1},b_t),\qquad c>0,
\label{eq:pairwise}
\end{equation}

where $b$, $c$, model parameters, and threshold are fixed over the threat set.
SemStamp and k-SemStamp use adjacent region transitions; SimMark uses a soft
count derived from adjacent sentence-embedding similarity. For fixed $m$,
their offsets and positive normalizers fit Equation~\ref{eq:pairwise}.

The formulation excludes global attention, higher-order unit interactions,
scan/post-selection statistics, stochastic scores, and thresholds that adapt
to the candidate segmentation.

\subsection{Exact dynamic programs}

For independent spans, let $F(t,j,r)$ be the minimum score sum for segmenting
the first $j$ tokens into $t$ units while introducing $r$ non-original
internal boundaries. The recurrence adds a legal span $[i,j)$ and increments
$r$ when the internal endpoint $j$ is not in $B_0$. With maximum span length
$W$, the time complexity is $O(mkLW)$ after span scores are available.

Pairwise detectors require the previous span. Let
$F(t,h,i,r)$ be the minimum pair-score sum for the first $i$ tokens split
into $t$ units, with final span $[h,i)$ and $r$ introduced boundaries.
Adding $[i,j)$ gives

\begin{equation}
F(t+1,i,j,r+\delta(j))
=
\min_h\left[F(t,h,i,r)+q(h,i,j)\right],
\label{eq:dp}
\end{equation}

where $\delta(j)=\mathbf 1[j\notin B_0]$ for $j<L$ and $\delta(L)=0$.
Only legal gaps, legal spans, $t\leq m$, and $r\leq k$ are retained. The exact
worst score is

\begin{equation}
C_k(x)=b+\frac{1}{c}
\min_{\substack{h\\0\leq r\leq k}}F(m,h,L,r).
\label{eq:certificate}
\end{equation}

The pairwise time and predecessor-space bounds are $O(mkLW^2)$ and
$O(mkLW)$, respectively.

\paragraph{Soundness.}
Every path accepted by Equation~\ref{eq:dp} contains exactly $m$ legal spans,
and its count $r$ equals the number of introduced boundaries. Conversely,
every $P\in\mathcal A_k$ induces exactly one such path. The recurrence keeps
the minimum score among paths with the same sufficient state. By induction on
$t$, Equation~\ref{eq:certificate} is therefore the minimum over
$\mathcal A_k$. Thus

\begin{equation}
C_k(x)>\tau
\quad\Longrightarrow\quad
\forall P\in\mathcal A_k(x,P_0),\ D(x,P)>\tau.
\end{equation}

This guarantee is exact for the declared lattice and detector class; it says
nothing about partitions outside that contract.

\subsection{Calibration}

Because $\mathcal A_k$ contains $P_0$ and expands with $k$,

\begin{equation}
C_{k+1}(x)\leq C_k(x)\leq C_0(x).
\end{equation}

Using a fixed raw upper-tail threshold cannot increase null FPR, but can make
the certificate unnecessarily conservative. For a power comparison at a
common target FPR $\alpha$, we calibrate each complete $C_k$ on an independent
null split. For $n$ calibration scores, the strict threshold is the
$r$-th ascending order statistic,

\begin{equation}
r=\left\lceil(n+1)(1-\alpha)\right\rceil,
\qquad \tau_k=C_{k,(r)}.
\end{equation}

Ties remain on the non-detection side because the decision is
$C_k>\tau_k$. A separate test null estimates realized FPR with a one-sided
Clopper--Pearson upper bound. Positive certification is reported with a
one-sided lower bound. These quantities are not pooled across radii.

\section{Algorithm and Systems Validation}

\subsection{Oracle and synthetic calibration results}

\begin{table}[t]
\caption{Model-independent validation of the exact framework.}
\label{tab:algorithm}
\centering
\small
\begin{tabular}{p{0.25\linewidth}rp{0.43\linewidth}}
\toprule
Check & Scale & Result \\
\midrule
Independent-span DP & 7,691 cases & Exact agreement with brute force \\
Pairwise DP & 7,669 cases & Exact score, count, decision, and witness \\
Pairwise-to-additive reduction & 1,000 cases & Exact agreement \\
300-token pairwise case & 14,606,509 partitions &
39,977 states, 449,337 transitions, 0.8054 s \\
Synthetic calibration & 99,980 nulls &
All 80 family/replicate/radius FPR gates passed \\
\bottomrule
\end{tabular}
\end{table}

The pairwise total in Table~\ref{tab:algorithm} comprises 5,669 systematic
and 2,000 random tied-score instances. The independent-span total comprises
5,691 systematic and 2,000 random instances and checks 36,824 legal
partitions. All witness scores are independently recomputed; $k=0$ returns
only the original segmentation; illegal or missing inputs fail closed.

In the synthetic calibration study, separately calibrated continuous-null
FPR means ranged from 4.57\% to 5.18\% over $k=0,\ldots,3$. The largest
one-sided 95\% FPR upper bound among 80 checks was 7.465\%, below the frozen
8\% gate. Reusing the raw $\tau_0$ at $k\geq1$ yielded zero FPR in these
synthetic families, illustrating conservatism rather than a real-watermark
result.

\subsection{Adapters and compiled core}

The real-system adapter preserves a typed unit-list tape and reconstructs the
original units byte-for-byte at $k=0$. Public-code inspection showed that
SemStamp, k-SemStamp, and SimMark are first-order adjacent-span systems, not
independent-span systems. Frozen small-sample tests then matched official and
direct computations: SemStamp passed 258/258 unit, region, mask, hit, and
score checks on 43 eligible documents; SimMark passed 240/240 preprocessing,
cosine, soft-count, and score checks on 60 documents. The SimMark check used
an explicitly diagnostic $\gamma=0.5$ sentinel and is not a detector-health
result.

The Python sparse implementation exceeded a 60-second gate on the largest
frozen real-data shape. A C++17 grouped, layer-local core preserved all oracle
outputs and processed a shape with 14,169 atoms, 845 units, and $k=3$ in
18.35 and 18.47 seconds using about 165 MB. It represented
$39{,}353{,}742{,}686{,}409{,}367{,}909$ legal segmentations through
961,430 unique spans and 38,664,070 unique pair queries.

Materializing pair tables then failed a five-million-pair gate. A factorized
core instead caches one feature per span and computes pair values on demand.
It passed real-provider equivalence for 12 documents per system, but full
two-system score-blind projections were 22.44 GPU-hours for SemStamp and
71.81 GPU-hours for SimMark. Increasing SimMark batch size gave only
1.0027$\times$ speedup, and an exploratory global-reuse audit reduced work by
only 0.6748\%; that audit also violated its frozen row-scope contract and is
marked invalid.

\subsection{Detector reproduction and health}

\begin{table}[t]
\caption{Real-system terminal results. Certificate rates are intentionally
absent when detector health or readiness failed.}
\label{tab:systems}
\centering
\small
\begin{tabular}{p{0.20\linewidth}rrp{0.40\linewidth}}
\toprule
Run & Test FPR & Positive rate & Interpretation \\
\midrule
SimMark v9, $k=0$ & 27/400 & 89/1000 &
Published-asset baseline unhealthy; no robustness inference \\
SimMark v9, $k=1$ & 37/400 & 97/1000 &
Recalibrated exact score, but invalid as watermark evidence \\
SemStamp v10, $k=0$ & 33/400 & 974/999 &
TPR passed; FPR CP95 upper 10.88\% failed the 8\% gate \\
SemStamp v17, $k=0$ & 47/1000 & 954/1000 &
Detector-health Go; CP95 bounds 5.95\% and 94.16\% \\
SemStamp v18, $k\geq1$ & -- & -- &
Readiness Stop before score or certificate computation \\
\bottomrule
\end{tabular}
\end{table}

SimMark v9 completed exact scores for 200 fit and 1,800
calibration/test/positive documents over $k=0,\ldots,3$, but the official
histogram produced $\gamma=0$ and clean-positive detection was only 8.9\%.
Subsequent diagnosis reproduced all 2,000 official/direct document scores and
localized the discrepancy to Instructor prompt-pooling compatibility. A
no-prompt-pooling diagnostic restored 99.0\% positive detection on the frozen
split, but it was not preregistered as a new detector and did not retroactively
validate v9.

SemStamp v10 detected 974/999 positives, but its test FPR was 8.25\% and the
one-sided 95\% upper bound was 10.88\%, above the frozen 8\% gate.
Across 100 deterministic rotated splits, this split was near the high end of
the observed FPR distribution. Rather than choose a favorable split, we
constructed fresh, identity-separated C4 cohorts and a new positive set.

The positive-generation sequence itself exposed reproducibility failures:
an initial smoke produced two prompt-only samples; one runner stalled because
its fallback was unreachable; a bounded replacement broke a strict-prefix
contract; and an independent audit later found five multi-sentence
prompt-only outputs among 1,000 claimed positives. A disclosed post-hoc
validation rescue retained 995 valid outputs and generated five replacements
from the fixed queue. This history is not rewritten as preregistered success.

With the resulting 1,000 positives and two independent 1,000-document null
cohorts, SemStamp v17 passed detector health. The threshold was the 951st
calibration order statistic, $\tau_0=1.6667$. Test FPR was 47/1000
(CP95 upper 5.95\%), positive TPR was 954/1000 (CP95 lower 94.16\%), and two
3,000-row replays were byte-identical. Official and direct scores agreed
exactly on 6,000 evaluations.

\subsection{Final exact-certificate readiness}

SemStamp v18 froze the same 3,000 identities, native tokenizer legality,
$k=0,\ldots,3$, a 50-million-pair per-document cap, and a 72 GPU-hour budget
for two complete replays. A score-blind pass reconstructed every workload
without loading the detector.

\begin{table}[t]
\caption{Frozen SemStamp v18 readiness gates.}
\label{tab:v18}
\centering
\begin{tabular}{lrrl}
\toprule
Gate & Limit & Observed maximum & Result \\
\midrule
Pruned pairs per document & 50,000,000 & 78,642,945 & Stop \\
All-allowed pairs per document & 50,000,000 & 83,540,900 & Stop \\
Two-replay budget & 72 h & 86.6501 h & Stop \\
Rows reconstructed & 3,000 & 3,000 & Pass \\
\bottomrule
\end{tabular}
\end{table}

The automatic gate left formal B1 as \texttt{STOPPED / Not started}.
Detector, model, region, score, threshold, generation, and certificate calls
were all zero. No threshold $\tau_k$, witness, test FPR, positive certified
rate, or non-vacuity statistic exists for v18. The result is a capacity and
budget limit of the validated exact implementation, not evidence that
SemStamp is or is not robust to resegmentation.

\section{Discussion}

\subsection{What the unified evidence establishes}

The SWAN observation shows that segmentation is not merely preprocessing:
it can determine the final decision. The exact framework then states what a
stronger robustness claim would require: enumerate a declared finite
boundary lattice implicitly, minimize the complete detector statistic, and
calibrate that statistic on independent null data.

The algorithmic result is positive. For independent-span and first-order
pairwise detectors, exhaustive search over all legal partitions can be
replaced by exact dynamic programming without weakening the threat set. The
systems result is mixed. Detector semantics, data identity, rendering,
feature extraction, calibration, and resource limits all matter before a
certified rate can be interpreted. On our final SemStamp assets, detector
health was adequate but the validated exact implementation exceeded its
predeclared capacity and runtime gates.

\subsection{Why the two streams do not substitute for each other}

The certificate threat set fixes token identity, token order, and unit count.
The strongest SWAN behavior result changes unit count. Conversely, the
SemStamp and SimMark detectors do not use SWAN's AMR parser or fixed AMR bank.
The common contribution is the segmentation-interface framing; causal and
statistical claims remain system-specific.

\subsection{Implications for watermark evaluation}

Sentence-level watermark reports should specify at least:

\begin{enumerate}
  \item the segmentation and rendering contract;
  \item whether unit count may change under the threat model;
  \item whether the detector statistic decomposes locally;
  \item whether thresholds are raw or calibrated after robustification;
  \item clean TPR and independent realized FPR before robust conclusions;
  \item the coverage and computational cost of any claimed certificate.
\end{enumerate}

Reporting only a successful attack or only a dynamic program is insufficient.
The former may collapse statistical power; the latter may be exact but
computationally or statistically vacuous.

\section{Limitations}

\begin{itemize}
  \item \textbf{SWAN sample scope.} The A--B--A cohort contains 30 English
  texts of 17--46 words and five original units. The normal-length,
  count-neutral endpoint was not run.
  \item \textbf{SWAN quality and calibration.} The short-text result lacks
  condition-specific FPR. The normal-length routes did not produce the
  independently human-reviewed cohort required by protocol.
  \item \textbf{Threat-set scope.} The certificate fixes tokens, order, and
  unit count. It does not certify arbitrary paraphrase, punctuation
  realizability, variable $N$, or semantic equivalence.
  \item \textbf{Detector class.} Exactness applies to independent-span or
  first-order adjacent-span scores with fixed affine aggregation. It does not
  cover global or higher-order detectors.
  \item \textbf{Real certificate endpoint.} SimMark's completed run had an
  unhealthy baseline. SemStamp's final healthy run stopped at readiness.
  No defensible real certified rate is reported.
  \item \textbf{Assets and environments.} Results depend on frozen public
  revisions, CUDA behavior, tokenizers, and generated SemStamp positives.
  The v16 positive set is a disclosed post-hoc validation rescue.
  \item \textbf{Multiplicity.} Radius-specific intervals are not presented as
  a simultaneous confidence statement over all radii.
\end{itemize}

\subsubsection*{Broader Impact Statement}

This work is intended to improve the evaluation and engineering of text
watermarks by making segmentation assumptions explicit. The same analysis
could inform attempts to reduce detectability. We limit that risk by
reporting aggregate frozen results, a bounded threat model, and no live key,
credential, sealed text, or service-targeting procedure. Exact certificates,
when available, can support defenders as well as auditors.

\section{Conclusion}

Sentence boundaries are part of the operational contract of sentence-level
watermarks. SWAN's decision changed reversibly under a short-text
segmentation intervention, but the available evidence does not isolate an
AMR-specific or normal-length count-neutral mechanism. Exact dynamic
programming can certify all segmentations in a fixed-token,
fixed-unit-count relocation ball for local detector classes, and our
implementation passes exhaustive correctness and calibration checks. The
final real-system result is nevertheless a Stop: SemStamp detector health
passed, while exact-certificate capacity and budget did not. This combination
of behavioral evidence, exact formulation, and audited non-completion defines
both the contribution and its boundary.

\bibliography{references}
\bibliographystyle{tmlr}

\appendix

\section{SWAN Exploration and Diagnostic Results}

\paragraph{Exploration history.}
Plans A--E explored ordinary paraphrasing, AMR node and edge edits, modifier
and argument operations, reentrancy changes, sentence restructuring, and
feedback-assisted candidate selection. Increasingly strong graph edits traded
off detectability against factual preservation. Plan F corrected denominators
to quality-gated valid attacks. Plans G--J isolated sentence boundaries,
tested no-feedback selection, compared SIRA-adapted, and added unmatched KGW
and SynthID-Text controls. These studies motivate the interface audit but are
not pooled as one confirmatory endpoint.

\paragraph{Artificial-graph diagnostic.}
Across 108 accepted injection targets, median target S2Match F1 was 0.8197
for original candidates, 0.4502 for artificial compositions, and 0.5200 for
saved merged-text reparses. Target-green counts were 108/108, 0/108, and
8/108. Artificial composition adds a root and connector edges and is not a
natural-text counterfactual.

\paragraph{Stopped defenses and range controls.}
Unmatched boundary-reorganization controls recorded valid attack rates of
0/30 for KGW and 1/30 for SynthID-Text. Three local SWAN defenses failed
their development gates: multi-instance rules recovered 0/15 across three
representations; template-tail detection reached 4/15 on original graphs and
2/15 on reparses, below a joint 8/15 gate; relation-motif recall at five was
41/56 and 33/56, each one below its continuation gate. A direct
resegmentation mitigation changed valid escapes from 23/30 to 19/30
($p=.125$) while human-null false positives changed from 3/250 to 5/250.
No effective segmentation-invariant SWAN defense is claimed.

\section{Certificate Engineering Timeline}

\begin{table}[h]
\caption{Terminal certificate stages. ``Pass'' applies only to the stated
gate, not to an end-to-end robustness claim.}
\label{tab:timeline}
\centering
\scriptsize
\begin{tabular}{p{0.10\linewidth}p{0.31\linewidth}p{0.50\linewidth}}
\toprule
Stage & Purpose & Terminal evidence \\
\midrule
WP1 & Additive exact DP & 7,691 oracle cases passed \\
WP1b & Pairwise edge-state DP & 7,669 oracle cases and 1,000 reductions passed \\
WP1c & Certificate calibration & 99,980 synthetic nulls; 80/80 FPR gates passed \\
WP2a & Canonical adapter & 500 model-free instances and $k=0$ restoration passed \\
v1--v2 & Public assets and lossless adapter &
Official/direct checks passed; frozen eligibility contracts failed \\
v3 & Typed structural separator & Eligibility recovered; Python core timed out \\
v4 & Compiled sparse core & Algorithm-Ready on the largest frozen shape \\
v5 & External pair table & Stopped at 5,000,001 pairs before model loading \\
v6 & Factorized features & Provider equivalence passed; SimMark projection 71.81 h \\
v7 & Batch scaling & Best speedup 1.0027$\times$; Stop \\
v8 & Global reuse & Invalid row scope; exploratory reuse only 0.6748\% \\
v9 & SimMark long run & Exact scores completed; detector baseline unhealthy \\
v10 & SemStamp health-first & TPR passed; independent FPR bound failed \\
v11 & Reproduction diagnosis &
Official/direct equivalence; SimMark pooling and SemStamp split instability \\
v12 & Fresh identities & 12,612 legal identities; smoke 18/20, then Stop \\
v13 & Replacement-aware generation & Stalled at attempt 372 after 84,981 s \\
v14 & Bounded fallback & Strict-prefix divergence at attempt 2; Stop \\
v15 & Causal-prefix generation & 1,000 claimed; audit found five prompt-only rows \\
v16 & Validation rescue & Disclosed 995+5 cohort; identity audit Go \\
v17 & SemStamp detector health & 4.7\% FPR, 95.4\% TPR; Go \\
v18 & Exact-certificate readiness & 78.6M pairs and 86.65 h; Stop before scoring \\
\bottomrule
\end{tabular}
\end{table}

No stopped version is retroactively relabeled as successful. The v9 scores
remain a detector-reproduction diagnostic; v10 remains a failed health gate;
v16 remains post-hoc; and v18 contains no certificate output.

\section{Additional SimMark Diagnostics}

The v9 radius-specific diagnostic rates were:

\begin{center}
\begin{tabular}{rrr}
\toprule
$k$ & Independent-test FPR & Positive rate \\
\midrule
0 & 27/400 & 89/1000 \\
1 & 37/400 & 97/1000 \\
2 & 36/400 & 64/1000 \\
3 & 36/400 & 34/1000 \\
\bottomrule
\end{tabular}
\end{center}

Each radius used its own calibrated threshold, so the positive point estimate
need not be monotone even though every document satisfies
$C_0\geq C_1\geq C_2\geq C_3$. The table is not a certified-robustness result
because the $k=0$ detector failed before-attack health.

\clearpage
\section{Unestimated Endpoints}

\begin{center}
\textbf{Table 8: Important quantities that must not be read as zero.}

\centering
\small
\begin{tabular}{p{0.26\linewidth}p{0.62\linewidth}}
\toprule
Endpoint & Status \\
\midrule
SWAN normal-length $N$-preserving effect & Not run \\
SWAN parser and candidate-graph mechanism & Not run \\
SWAN condition-specific FPR & Not run \\
SemStamp exact $C_k$ on final v18 cohorts & Not run \\
SemStamp post-certificate $\tau_k$, FPR, TPR & Not run \\
SemStamp real certified-positive rate & Not run \\
Healthy SimMark certified-positive rate & Not available \\
Cross-system certified robustness comparison & Not available \\
\bottomrule
\end{tabular}
\end{center}

\section{Reproducibility Boundary}

The internal evidence tree stores frozen protocols, code and model revisions,
identity manifests, row-level derived outputs, independent audits, terminal
states, and SHA256 chains. The SWAN reconciliation audit preserves 39,457
file hashes and 82 original-document copies. The SemStamp v17 health run uses
two byte-identical 3,000-row replays; the v18 workload and timing rows are
independently recomputed without loading the detector.

A public artifact must exclude credentials, sealed source or generated text,
and third-party payloads without redistribution permission. It can include
the claim registry, aggregate tables, exact-DP and audit code, frozen
revisions, non-sensitive manifests, and synthetic tests. Local paths are
reproducibility indices, not publication URLs.

\end{document}
